\section{Adaptive Learning Engine}

\label{sec:adaptive}

\subsection{Learner Model}

We model each learner's knowledge state using an augmented Bayesian Knowledge Tracing (BKT) framework.
For learner $u$ and objective $o_i$, the mastery state $M_{u,i}^{(t)} \in \{0, 1\}$ at interaction step $t$ follows:

\begin{equation}
    P(M_{u,i}^{(t)} = 1 \mid M_{u,i}^{(t-1)} = 0) = p_T(o_i, u)
\end{equation}

\begin{equation}
    P(M_{u,i}^{(t)} = 0 \mid M_{u,i}^{(t-1)} = 1) = p_F(o_i, u)
\end{equation}

\noindent where $p_T$ is the transition probability (learning rate) and $p_F$ is the forgetting probability.
Unlike standard BKT, we parameterize both $p_T$ and $p_F$ as functions of learner features and objective characteristics:

\begin{equation}
    p_T(o_i, u) = \sigma\left(\mathbf{w}_T^\top [\mathbf{e}_{o_i}; \mathbf{h}_u; \mathbf{f}_{u,i}]\right)
\end{equation}

\noindent where $\sigma$ is the sigmoid function, $\mathbf{e}_{o_i}$ is the objective embedding from the CCO graph (obtained via graph neural network encoding), $\mathbf{h}_u$ is the learner's latent profile vector, and $\mathbf{f}_{u,i}$ captures interaction features (time spent, attempt count, recency).

\subsection{Content Selection Policy}

Given the learner model, ZooEd selects the next learning activity by optimizing an acquisition function that balances learning efficiency, engagement, and ecological relevance:

\begin{equation}
    o^* = \arg\max_{o_i \in \mathcal{F}_u} \left[\alpha \cdot \text{IG}(o_i, u) + \beta \cdot \text{Eng}(o_i, u) + \gamma \cdot \text{Eco}(o_i, \mathbf{s}_t)\right]
\end{equation}

\noindent where $\mathcal{F}_u$ is the set of feasible objectives (prerequisites satisfied), $\text{IG}$ is the expected information gain for the learner model, $\text{Eng}$ is the predicted engagement score, $\text{Eco}$ is the ecological relevance score given current sensor state $\mathbf{s}_t$, and $\alpha, \beta, \gamma$ are tunable weights.

The information gain component is computed as the expected reduction in entropy of the learner's knowledge state:

\begin{equation}
    \text{IG}(o_i, u) = H(M_{u,i}) - \mathbb{E}_{Y}\left[H(M_{u,i} \mid Y)\right]
\end{equation}

\noindent where $Y$ is the predicted assessment outcome.

\subsection{Spaced Repetition with Ecological Context}
\label{sec:srec}

Standard spaced repetition algorithms schedule reviews based on forgetting curves estimated from past performance \citep{settles2016trainable}.
Our SREC algorithm extends this by incorporating ecological triggers:

\begin{equation}
    r_i(t) = \max\left(r_i^{\text{base}}(t),\ r_i^{\text{eco}}(t)\right)
\end{equation}

\noindent where $r_i^{\text{base}}(t)$ is the base review urgency from the forgetting model and $r_i^{\text{eco}}(t)$ is an ecological trigger score.
The ecological trigger activates when sensor data from the Zoo network detects conditions relevant to a learned objective:

\begin{equation}
    r_i^{\text{eco}}(t) = \text{sim}(\mathbf{e}_{o_i}, \phi(\mathbf{s}_t)) \cdot \mathbb{1}[\Delta t_{u,i} > \tau_{\min}]
\end{equation}

\noindent where $\phi(\mathbf{s}_t)$ projects the sensor state into the objective embedding space and $\Delta t_{u,i}$ is the time since the learner last engaged with objective $o_i$, gated by a minimum interval $\tau_{\min}$ to prevent notification fatigue.

For example, when camera traps in a learner's region detect a species they recently studied, SREC triggers a contextual review activity incorporating the actual detection image, reinforcing the connection between learned knowledge and field observation.

\subsection{Multi-Language Content Generation}

ZooEd generates learning content in 47 languages using a pipeline architecture:

\begin{enumerate}[leftmargin=*]
    \item \textbf{Canonical Content:} Subject matter experts author core content in English with structured annotations marking cultural adaptability points.
    \item \textbf{Translation:} A multilingual transformer model fine-tuned on ecological text produces initial translations, with ecological terminology mapped through a curated glossary of 12,400 terms across all supported languages.
    \item \textbf{Cultural Adaptation:} Local experts review and adapt examples, analogies, and cultural references. This review is incentivized through Zoo token rewards.
    \item \textbf{Quality Assurance:} Back-translation verification and automated readability scoring ensure translations maintain pedagogical clarity.
\end{enumerate}

The content generation model is formalized as:

\begin{equation}
    \mathbf{y}_{\ell} = G(\mathbf{x}_{\text{en}}, \ell, \mathbf{c}_{\text{region}}, \mathbf{m}_{\text{glossary}})
\end{equation}

\noindent where $G$ is the generation model, $\mathbf{x}_{\text{en}}$ is English source content, $\ell$ is the target language code, $\mathbf{c}_{\text{region}}$ is regional context features, and $\mathbf{m}_{\text{glossary}}$ is the ecological terminology mapping.

% ============================================================
% 5. Assessment Framework
% ============================================================
